% from ega0-19-fr.tex
\oldpage[0\textsubscript{IV-1}]{91}

Lorsque $\mathbf S_C^n(\mathfrak J/\mathfrak J^2)$ est muni de la topologie
définie dans (19.5.1), les homomorphismes $\varphi_n$ sont continus, en vertu
de la propriété universelle (19.5.1) de
$\mathbf S_C^*(\mathfrak J/\mathfrak J^2)$ appliquée à la $C$-algèbre
topologique $E=\gr_{\mathfrak J/\mathfrak J^n}(B/\mathfrak J^n)$ munie de la
topologie produit de celles des $\mathfrak J^k/\mathfrak J^{k+1}$, et à
l'injection canonique $\mathfrak J/\mathfrak J^2\to E$. Notons qu'ici la
topologie de $\mathfrak J/\mathfrak J^2$ est moins fine que la topologie
déduite de celle de $C$ (cette dernière topologie sur
$\mathfrak J/\mathfrak J^2$ étant aussi la topologie déduite de celle de $B$).
\end{env}

\begin{theorem}[19.5.3]
\label{0.19.5.3-fr}
Soient $A$ un anneau topologique, $B$ une $A$-algèbre topologique,
$\mathfrak J$ un idéal de $B$, $C=B/\mathfrak J$ la $A$-algèbre topologique
quotient. On suppose que la $A$-algèbre discrète $C$ est formellement lisse.
Alors~:
\begin{enumerate}
  \item[\rm(i)] Si $B$ est une $A$-algèbre formellement lisse,
  $\mathfrak J/\mathfrak J^2$ est un $C$-module topologique formellement
  projectif (19.2.1).
  \item[\rm(ii)] Si $B$ est une $A$-algèbre formellement lisse et un anneau
  préadmissible ($\mathbf 0_{\rm I}$, 7.1.2), les homomorphismes $\varphi_n$
  (19.5.2) sont des bimorphismes formels (19.1.2).
  \item[\rm(iii)] Réciproquement, supposons que $B$ soit préadmissible, que la
  suite $(\mathfrak J^n)$ tende vers $0$ dans $B$, que
  $\mathfrak J/\mathfrak J^2$ soit un $C$-module formellement projectif et que
  les $\varphi_n$ soient des bimorphismes formels. Alors $B$ est une
  $A$-algèbre formellement lisse.
\end{enumerate}
\end{theorem}

La démonstration de ce théorème est longue et encombrée de détails techniques~;
aussi commencerons-nous par en démontrer un corollaire plus simple, où
apparaissent plus nettement les idées directrices~; ce corollaire est d'ailleurs
le cas particulier du théorème (19.5.3) qui sera le plus fréquemment employé
par la suite.

\begin{corollary}[19.5.4]
\label{0.19.5.4-fr}
Soient $A$ un anneau topologique, $B$ une $A$-algèbre topologique,
$\mathfrak J$ un idéal de $B$ tel que la topologie de $B$ soit la topologie
$\mathfrak J$-préadique. On suppose que la $A$-algèbre discrète
$C=B/\mathfrak J$ est formellement lisse. Alors les trois conditions suivantes
sont équivalentes~:
\begin{enumerate}
  \item[a)] $B$ est une $A$-algèbre formellement lisse.
  \item[b)] $\mathfrak J/\mathfrak J^2$ est un $C$-module projectif et
  l'homomorphisme canonique
  \[
    \varphi:\mathbf S_C^*(\mathfrak J/\mathfrak J^2)
    \longrightarrow\gr_{\mathfrak J}^*(B)
  \]
  est bijectif.
  \item[c)] Le séparé complété $\widehat B$ de $B$ est une
  $\widehat A$-algèbre topologique isomorphe à une $\widehat A$-algèbre de la
  forme $\widehat{B'}$, où $B'=\mathbf S_C^*(V)$, $V$ étant un $C$-module
  projectif et $B'$ étant muni de la topologie ${B'}^+$-préadique, où
  ${B'}^+$ est l'idéal d'augmentation.
\end{enumerate}
\end{corollary}

\begin{env}[19.5.4.1]
\label{0.19.5.4.1-fr}
Prouvons d'abord que $a)$ entraîne que $\mathfrak J/\mathfrak J^2$ est un
$C$-module projectif. Soient donc $P$ et $Q$ deux $C$-modules,
$u:P\to Q$ un $C$-homomorphisme surjectif, et
$v:\mathfrak J/\mathfrak J^2\to Q$ un $C$-homomorphisme. L'anneau
$B/\mathfrak J^2$ étant considéré comme une $B$-extension de $C$ par
$\mathfrak J/\mathfrak J^2$, on en déduit par $v$ une $B$-extension
$E=(B/\mathfrak J^2)\oplus_{(\mathfrak J/\mathfrak J^2)}Q$ de $C$ par $Q$
(18.2.8). Comme $C$ est une $A$-algèbre discrète formellement lisse,
l'extension $E$ est $A$-\emph{triviale} (19.4.4) et peut donc s'identifier à
$\mathrm D_C(Q)$ (18.2.3). L'homomorphisme surjectif $u:P\to Q$ définit alors
canoniquement un $A$-homomorphisme surjectif
$\psi:\mathrm D_C(P)\to\mathrm D_C(Q)$ (18.2.9) dont le noyau est un idéal
$\mathfrak L$ de l'extension $E'=\mathrm D_C(P)$, contenu dans $P$, et
\emph{a fortiori} de carré nul. Soit $f:B\to E=E'/\mathfrak L$
l'homomorphisme définissant la structure de $B$-algèbre de $E$~: comme
$\mathfrak L$ est de carré nul, et que $B$ est une $A$-algèbre formellement
lisse, $f$ se factorise en $B\xrightarrow{g}E'\to E'/\mathfrak L$, où $g$ est
un $A$-homomorphisme continu.

